%% kapa: exam=1 round=26 year=2015 period=2 form=A subject=accounting subject_ko=회계학 seq=26 number=66
%% source: Q-Net 감정평가사 기출문제 (cst00309) artlSeq=5138001 file=extracted/1008/5138001/제26회_감정평가사 1차_기출/제26회_감정평가사 1차_기출/2교시_감평_A형.pdf
%% conversion: kordoc 4.24.51 -> tools/gen1.py (auto); review=proofread-sonnet; flags=none
\begin{question}{66}
(주)감평은 제품 구입 후 1년 이내에 발생하는 제품의 결함에 대하여 제품보증을 실시하고 있다. 20×3년에 판매된 제품에 대하여 중요하지 않은 결함이 발생한다면 ￦50,000의 수리비용이 발생하고, 치명적인 결함이 발생하면 ￦200,000의 수리비용이 발생한다. 과거경험률에 따르면 70\%는 결함이 없으며, 20\%는 중요하지 않은 결함이 발생하며, 10\%는 치명적인 결함이 발생한다고 할 때 20×3년 말에 제품보증충당부채로 인식할 금액은? (단, 20×3년 말까지 발생한 수리비용은 없다.)
\begin{choices}
  \item ￦10,000
  \item ￦20,000
  \item ￦30,000
  \item ￦200,000
  \item ￦250,000
\end{choices}
\end{question}
\answer{3}
